
%\documentclass[./main.tex(path of main file)]{subfiles}
\documentclass[10pt]{article}
\usepackage{amsmath}
\DeclareMathOperator*{\argmin}{arg\,min} % thin space, limits underneath in displays
\DeclareMathOperator*{\argmax}{arg\,max} % thin space, limits underneath in displays
\newtheorem{thm}{Theorem}
\usepackage{amssymb}
\usepackage{amsfonts}
\usepackage{mathrsfs}
\usepackage{bm}
\usepackage{indentfirst}
\setlength{\parindent}{0em}
\usepackage[margin=1in]{geometry}
\usepackage{graphicx}
\usepackage{setspace}
\doublespacing
\usepackage[flushleft]{threeparttable}
\usepackage{booktabs,caption}
\usepackage{float}
\usepackage[sort,comma]{natbib}
\usepackage[hidelinks]{hyperref}
\usepackage{booktabs}
\usepackage{multirow}

\usepackage{import}
\usepackage{xifthen}
\usepackage{pdfpages}
\usepackage{transparent}
\usepackage{subfiles}
\usepackage{blindtext}
\usepackage{glossaries}
\usepackage{pdflscape}
\usepackage{adjustbox}
\usepackage{rotating}
\usepackage{eurosym,geometry,ulem,subcaption,caption,color,sectsty,comment,footmisc,pdflscape,array}

% Format definition list
\newenvironment{mylist}[1]%
  {\begin{list}{}{\settowidth{\labelwidth}{\textbf{#1}}
   \setlength{\leftmargin}{\labelwidth}
   \addtolength{\leftmargin}{\labelsep}
   \renewcommand{\makelabel}[1]{\textbf{\hfill##1}}}}%
  {\end{list}}



\newcommand{\incfig}[1]{%
\def\svgwidth{\columnwidth}
\import{./figures/}{#1.pdf_tex}
}




\title{}
\author{}
\date{}


\begin{document}
%\maketitle

\tableofcontents
\newpage


\section{The Permanent-Income Hypothesis (Consumption under Certainty)}


\subsection{Permanent Income}

Consider a simple utility function.
The Lagrangian for the maximization problem is
\begin{equation}
L = \sum\limits_{i = 1} ^T u(C_{t}) + \lambda
\left( 
		A_0 + \sum\limits_{t = 1} ^T Y_{t} - \sum\limits_{t = 1} ^T C_{t}		
\right) .
\end{equation}
The FOC for $ C_{t} $ is 
\begin{equation*}
u'(C_{t}) = \lambda
\end{equation*}

Since the above FOC holds every period, the marginal utility of consumption is constant.
And since the level of consumption uniquely determines its marginal utility, this meant that
consumption must be constant, i.e., $ C_1 = C_2 = ... = C_{T} $.
Substituting this fact into the BC yields,
\begin{equation}
C_{t} = \frac{1}{T}
\left( 
		A_0 + \sum\limits_{t = 1} ^T Y_{t}	
\right) , \text{ for all } t. \label{eqn:permanent_income_C}
\end{equation}

The term in the parentheses is the individual's total lifetime resources. 
Thus equation \eqref{eqn:permanent_income_C} states that 
{\textbf {the individual divides his or her lifetime resources equally among each period of life}}.
This analysis implies that the individual's consumption in a given period is determined
{\underline {NOT}} by income that period, but by income over his or her entire life.
In the terminology of Friedman (1957), the RHS of equation \eqref{eqn:permanent_income_C} is
{\underline {permanent income}} (it is an averaged income), and the difference between current
and permanent income is {\underline {transitory income}}.
{\textbf {
Equation \eqref{eqn:permanent_income_C} implies that consumption is determined by permanent income.
}}

The permanent-income hypothesis assumes that individuals can borrow at the same interest rate at
which they can save as long as they eventually repay their loans.



\subsection{Savings}
\begin{align}
S_{t} &= Y_{t} - C_{t}\\
&= 
\left( 
		Y_{t} - \frac{1}{T} \sum\limits_{\tau = 1} ^T Y_{\tau}	
\right)  
- \frac{1}{T} A_0 \label{eqn:permanent_income_S}
\end{align}

Savings is high when current income ($ Y_{t} $) is high relative to its average 
($ \frac{1}{T} \sum\limits Y_{\tau}	 $)--that is, when transitory income is high.
Thus {\underline {the individual uses saving and borrowing to smooth the path of consumption.}}
This is the key idea of the permanent-income hypothesis of Modigliani and Brumberg (1954) and
Friedman (1957).

At a more general level, saving is future consumption.




\subsection{Conclusion}
The permanent-income hypothesis provides appealing explanations of many important features of 
consumption. For example, it explains why temporary tax cuts appear to have much smaller
effects than permanent ones, and it accounts for many features of the relationship between current
income and consumption.

Yet there are also important features of consumption that appear inconsistent with the 
permanent-income hypothesis. For example, both macroeconomic and microeconomic evidence suggest
that consumption often responds to predictable changes in income. And simple models of consumer
optimization cannot account for the equity premium.

Indeed, the permanent-income hypothesis fails to explain some central features of consumption
behavior. One of the hypothesis's key predictions is that there should be no relation between the
expected growth of an individual's income over his or her lifetime and the expected growth of
his or her consumption: consumption growth is determined by the real interest rate and the discount
rate, not by the time pattern of income.

Carroll and Summers (1991) present extensive evidence that this prediction is incorrect. For example,
individuals in countries where income growth is high typically have high rates of consumption
growth over their lifetime, and individuals in slowly growing countries typically have low rates
of consumption growth.

More generally, most households have little wealth, and their consumption approximately tracks
their income. As a result, their current income has a large role in determining their consumption.




\section{The Random Walk Hypothesis (Consumption under Uncertainty)}





\section{Precautionary Saving}
The combination of a positive third derivative of the utility function ($ U'(C) $ is convex to the
origin) and uncertainty about future income reduces current consumption, and thus raises saving.
This saving is known as {\underline {precautionary saving}} (Leland, 1968).




\section{Liquidity Constraints}

The permanent-income hypothesis assumes that individuals can borrow at the same interest rate at
which they can save as long as they eventually repay their loans.
Yet the interest rates that households pay on credit card debt, automobile loans, and other
borrowing are often much higher than the rates they obtain on their savings. In addition, some 
individuals are unable to borrow more at any interest rate.

{\textbf {Liquidity constraints can raise saving in two ways}}. First, and most obviously, whenever a
liquidity constraint is binding, it causes the individual to consume less than he or she otherwise
would. Second, even if the constraints are not currently binding, the fact that they may bind in the
future reduces current consumption. Suppose, for example, there is some chance of low income
in the next period. If there are no liquidity constraints and income turns out to be low, the 
individual can borrow to avoid a sharp fall in consumption. If there are liquidity constraints, 
however, a large fall in income causes a large fall in consumption unless the individual
has savings. Thus liquidity constraints cause individuals to save as insurance against the
effects of future falls in income.






\bibliographystyle{plainnat}
\bibliography{my_bib}

\end{document}

